% ============================================================
% Paper 90 (DE): Die Bunyakovsky-Vermutung und irreduzible
%                Polynome, die Primzahlen darstellen
% Autor: Michael Fuhrmann
% Specialist Mathematics Project, Build 168
% Datum: 2026-03-12
% ============================================================
\documentclass[12pt,a4paper]{amsart}

\usepackage{amsmath,amssymb,amsthm,mathtools,enumitem,booktabs,array,hyperref}
\usepackage[utf8]{inputenc}
\usepackage[T1]{fontenc}
\usepackage{geometry}
\geometry{margin=2.5cm}

% ---- Theorem-Umgebungen (Deutsch) ----
\newtheorem{theorem}{Theorem}[section]
\newtheorem{lemma}[theorem]{Lemma}
\newtheorem{corollary}[theorem]{Korollar}
\newtheorem{proposition}[theorem]{Proposition}
\newtheorem{satz}[theorem]{Satz}
\newtheorem{korollar}[theorem]{Korollar}
\theoremstyle{definition}
\newtheorem{definition}[theorem]{Definition}
\newtheorem{definition_de}[theorem]{Definition}
\newtheorem{example}[theorem]{Beispiel}
\theoremstyle{remark}
\newtheorem{remark}[theorem]{Bemerkung}
\newtheorem{bemerkung}[theorem]{Bemerkung}
\newtheorem{conjecture}[theorem]{Vermutung}
\newtheorem{vermutung}[theorem]{Vermutung}

% ---- Operatoren ----
\newcommand{\N}{\mathbb{N}}
\newcommand{\Z}{\mathbb{Z}}
\newcommand{\Q}{\mathbb{Q}}
\newcommand{\R}{\mathbb{R}}

% ============================================================
\begin{document}

\title{Die Bunyakovsky-Vermutung und irreduzible\\
Polynome, die Primzahlen darstellen}

\author{Michael Fuhrmann}
\address{Specialist Mathematics Project, Build 168}
\email{specialist-maths@project.local}
\date{2026-03-12}

\begin{abstract}
Im Jahr 1857 vermutete Viktor Bunyakovsky, dass jedes irreduzible Polynom
$f \in \Z[x]$ mit positivem führenden Koeffizienten, das eine einfache
notwendige Bedingung erfüllt --- dass die Werte $f(n)$ für $n \in \Z$ keinen
gemeinsamen Teiler $> 1$ besitzen --- unendlich viele Primzahlwerte annimmt.
Dies ist die natürliche Verallgemeinerung von Dirichlets Satz über Primzahlen
in arithmetischen Progressionen (der lineare Fall).
Für Polynome vom Grad $\geq 2$ ist die Vermutung vollständig offen.
Diese Arbeit präsentiert die genaue Formulierung, die notwendigen lokalen
Bedingungen, die Bateman--Horn-Vermutung über das Zählen von Darstellungen,
Schinzels Hypothese H (eine Verallgemeinerung auf simultane Polynomwerte),
explizite Beispiele (darunter das Problem, ob $n^2 + 1$ unendlich oft prim ist),
sowie die Verbindung zu Eulers glücklichen Zahlen.
\end{abstract}

\maketitle

\tableofcontents

% ============================================================
\section{Einleitung}
% ============================================================

Der grundlegendste Satz über Primzahlen in Folgen ist Dirichlets Satz (1837):

\begin{satz}[Dirichlet 1837]
Ist $\gcd(a, b) = 1$, so enthält die arithmetische Progression
$a, a+b, a+2b, \ldots$ unendlich viele Primzahlen.
\end{satz}

Das erzeugende Polynom dieser Progression ist $f(n) = a + bn$ --- linear,
irreduzibel über $\Q$, mit positivem führenden Koeffizienten und
$\gcd_{n \in \Z}\{f(n)\} = 1$.

Die Bunyakovsky-Vermutung fragt: Gilt dasselbe für Polynome höheren Grades?

\begin{vermutung}[Bunyakovsky 1857]
\label{verm:buny}
Sei $f \in \Z[x]$ ein irreduzibles Polynom vom Grad $d \geq 1$ mit
positivem führenden Koeffizienten.
Gilt $\gcd_{n \geq 0} f(n) = 1$ (keine feste Primzahl teilt alle Werte),
so nimmt $f(n)$ für unendlich viele $n \in \N$ Primzahlwerte an.
\end{vermutung}

\begin{bemerkung}
Die Bedingung $\gcd_{n \geq 0} f(n) = 1$ ist eine rein lokale Bedingung:
Es gibt keine Primzahl $p$ mit $p \mid f(n)$ für alle $n$.
\end{bemerkung}

Für $d = 1$ ist Vermutung~\ref{verm:buny} Dirichlets Satz (bewiesen).
Für $d \geq 2$ ist sie vollständig offen.

% ============================================================
\section{Die lokale Obstruktionsbedingung}
% ============================================================

\begin{definition_de}
Eine Primzahl $p$ ist ein \emph{fester Teiler} von $f$, wenn $p \mid f(n)$
für alle $n \in \Z$ gilt.
Der relevante \emph{Inhalt} ist
\[
  \Delta(f) = \gcd\{f(n) : n \in \Z\}.
\]
\end{definition_de}

\begin{proposition}
\label{prop:lokal}
Für $f \in \Z[x]$ vom Grad $d$ gilt:
\[
  \Delta(f) = \gcd\{f(0), f(1), \ldots, f(d)\}.
\]
\end{proposition}

\begin{proof}
Die Werte $f(0), \ldots, f(d)$ bestimmen $f$ eindeutig (Lagrange-Interpolation).
Jede Primzahl $p \mid \Delta(f)$ teilt $f(n)$ für alle ganzen Zahlen $n$
durch ein Periodizitäts- und Bewertungsargument.
\end{proof}

\begin{example}
$f(n) = n^2 + n + 2$: $\gcd(f(0), f(1), f(2)) = \gcd(2, 4, 8) = 2$.
Also ist $2$ ein fester Teiler: $f(n)$ nimmt nie einen ungeraden Primwert an.
\end{example}

\begin{example}
$f(n) = n^2 + 1$: $\Delta(f) = 1$.
Die ersten Primwerte: $f(1)=2$, $f(2)=5$, $f(4)=17$, $f(6)=37$, $f(10)=101$,
$f(14)=197$.
\end{example}

\begin{vermutung}[$n^2 + 1$ unendlich oft prim]
\label{verm:nsq}
Es gibt unendlich viele $n \in \N$, für die $n^2 + 1$ prim ist.
\end{vermutung}

Vermutung~\ref{verm:nsq} ist ein Spezialfall der Bunyakovsky-Vermutung
und vollständig offen.

% ============================================================
\section{Lokale Bedingungen: notwendig aber nicht hinreichend}
% ============================================================

\begin{satz}[Lokale Bedingungen sind notwendig]
Hat $f \in \Z[x]$ einen festen Primteiler $p$, so ist $p \mid f(n)$ für alle
$n$, und $f(n)$ kann höchstens einen Primzahlwert ($p$ selbst) annehmen.
\end{satz}

\begin{bemerkung}
Die Bunyakovsky-Vermutung behauptet, dass lokale Bedingungen auch
\emph{hinreichend} für die Unendlichkeit der Primzahlwerte sind.
Dies ist das tiefe, unbewiesene Kernstück.
\end{bemerkung}

\begin{example}[Lokale Bedingungen genügen nicht zur Reduktionsprüfung]
$f(n) = n^4 + 4$: Durch Sophie Germains Identität gilt
\[
  n^4 + 4 = (n^2 + 2n + 2)(n^2 - 2n + 2),
\]
also ist $f$ für alle $n > 1$ zusammengesetzt, obwohl $\Delta(f) = 1$.
Aber $f$ ist \emph{reduzibel} über $\Q$ --- die Irreduzibilitätsbedingung
der Bunyakovsky-Vermutung ist nicht erfüllt.
\end{example}

Dieses Beispiel zeigt, warum Irreduzibilität wesentlich ist.

% ============================================================
\section{Dirichlets Satz als linearer Spezialfall}
% ============================================================

\begin{satz}[Dirichlet 1837, vollständige Aussage]
Für $\gcd(a, b) = 1$ enthält $(a + bn)_{n \geq 0}$ unendlich viele Primzahlen.
Außerdem sind die Primzahlen gleichmäßig auf die zu $b$ teilerfremden
Restklassen verteilt:
\[
  \pi(x; b, a) \sim \frac{1}{\phi(b)} \cdot \frac{x}{\log x}
  \quad (x \to \infty).
\]
\end{satz}

\begin{proof}[Beweisidee]
Verwendet Dirichlet-$L$-Reihen $L(s, \chi) = \sum_{n=1}^\infty \chi(n) n^{-s}$
für Dirichlet-Charaktere $\chi$ modulo $b$ und zeigt $L(1, \chi) \neq 0$
für $\chi \neq \chi_0$.
\end{proof}

\begin{bemerkung}
Der Beweis von Dirichlets Satz nutzt die Gruppenstruktur von $(\Z/b\Z)^*$
in wesentlicher Weise.
Diese Struktur ist für Polynome höheren Grades nicht verfügbar --- daher
ist die Verallgemeinerung so schwierig.
\end{bemerkung}

% ============================================================
\section{Die Bateman--Horn-Vermutung}
% ============================================================

Bateman und Horn (1962) formulierten eine quantitative Version der
Bunyakovsky-Vermutung:

\begin{vermutung}[Bateman--Horn 1962]
\label{verm:BH}
Seien $f_1, \ldots, f_k \in \Z[x]$ verschiedene irreduzible Polynome mit
positiven führenden Koeffizienten und $\Delta(f_1 \cdots f_k) = 1$.
Die Anzahl der $n \leq x$, für die alle $f_1(n), \ldots, f_k(n)$ simultan
prim sind, ist asymptotisch
\[
  \frac{C(f_1, \ldots, f_k)}{d_1 \cdots d_k} \cdot \frac{x}{(\log x)^k},
\]
wobei $d_i = \deg f_i$ und die \emph{singuläre Reihe}
\[
  C(f_1, \ldots, f_k)
  = \prod_p \left(1 - \frac{1}{p}\right)^{-k}
    \left(1 - \frac{\omega(p)}{p}\right)
\]
mit $\omega(p)$ gleich der Anzahl der Lösungen von
$f_1(n) \cdots f_k(n) \equiv 0 \pmod{p}$.
\end{vermutung}

\begin{bemerkung}
Für $k=1$, $f_1(n) = n^2 + 1$ sagt Vermutung~\ref{verm:BH} voraus:
\[
  |\{n \leq x : n^2 + 1 \text{ prim}\}| \sim \frac{C}{2} \cdot \frac{x}{\log x},
\]
wobei $C = \prod_{p \equiv 1 (4)} \frac{p-1}{p-2} \cdot \prod_{p \equiv 3 (4)} \frac{p}{p-1}$.
\end{bemerkung}

% ============================================================
\section{Schinzels Hypothese H}
% ============================================================

\begin{vermutung}[Schinzel--Sierpiński Hypothese H, 1958]
\label{verm:H}
Seien $f_1, \ldots, f_k \in \Z[x]$ irreduzible Polynome mit positiven
führenden Koeffizienten. Hat das Produkt $f_1 \cdots f_k$ keinen festen
Primteiler, so gibt es unendlich viele $n \in \N$, für die
$f_1(n), \ldots, f_k(n)$ simultan prim sind.
\end{vermutung}

\begin{bemerkung}
Hypothese H impliziert:
\begin{itemize}[nosep]
  \item Die Bunyakovsky-Vermutung ($k=1$, $\deg f \geq 2$);
  \item Die Primzahlzwillings-Vermutung ($f_1(n) = n$, $f_2(n) = n+2$);
  \item Die Primkonstellation-Vermutung;
  \item Unendlich viele Sophie-Germain-Primzahlen.
\end{itemize}
Hypothese H gilt als mit aktuellen Methoden völlig unerreichbar.
\end{bemerkung}

% ============================================================
\section{Eulers Glückszahlen und primerzeugende Polynome}
% ============================================================

\begin{example}[Eulers primerzeugendes Polynom]
Das Polynom $f(n) = n^2 - n + 41$ liefert Primzahlen für $n = 0, 1, \ldots, 40$:
\[
  f(0) = 41,\quad f(1) = 41,\quad f(2) = 43,\quad \ldots,\quad f(40) = 1601,
\]
alle prim. Aber $f(41) = 41^2 = 1681$ ist zusammengesetzt.
\end{example}

\begin{definition_de}
Eine positive ganze Zahl $m$ heißt \emph{Euler-Glückszahl}, wenn
$n^2 - n + m$ für $n = 1, 2, \ldots, m-1$ stets prim ist.
\end{definition_de}

\begin{satz}[Klassifikation der Euler-Glückszahlen]
Die einzigen Euler-Glückszahlen sind $2, 3, 5, 11, 17, 41$.
\end{satz}

\begin{bemerkung}
Das Polynom $n^2 + n + 41$ hat ebenfalls bemerkenswerte primerzeugende
Eigenschaften.
Unter der Bunyakovsky-Vermutung stellt es unendlich viele Primzahlen dar.
Die außergewöhnlichen Eigenschaften dieses Polynoms hängen mit dem
imaginär-quadratischen Körper $\Q(\sqrt{-163})$ zusammen, der
Klassenzahl $1$ hat (Baker--Heegner--Stark).
\end{bemerkung}

% ============================================================
\section{Verbindung zur algebraischen Geometrie}
% ============================================================

\begin{bemerkung}
Ken Ono und Mitarbeiter haben Verbindungen zwischen der Bunyakovsky-Vermutung
und der Arithmetik algebraischer Kurven untersucht.
Das Problem, ob $f(n)$ prim ist, hängt mit rationalen Punkten auf gewissen
algebraischen Varietäten zusammen.
\end{bemerkung}

\begin{satz}[Ono, partielles Ergebnis]
Sei $f(n) = n^2 + an + b$ irreduzibel über $\Q$ mit $\Delta(f) = 1$.
Hat der zugehörige imaginär-quadratische Körper $\Q(\sqrt{a^2-4b})$
die Klassenzahl $1$, so zerfallen alle Primwerte von $f$ in diesem Körper.
\end{satz}

\begin{bemerkung}
Die imaginär-quadratischen Körper der Klassenzahl $1$ sind
$\Q(\sqrt{-d})$ für $d \in \{1, 2, 3, 7, 11, 19, 43, 67, 163\}$
(Baker--Heegner--Stark).
Das Polynom $n^2 + n + 41$ entspricht $\Q(\sqrt{-163})$, was seine
außergewöhnliche Primzahlproduktion erklärt.
\end{bemerkung}

% ============================================================
\section{Das polynomiale Goldbach-Problem}
% ============================================================

\begin{vermutung}[Polynomiales Goldbach]
\label{verm:polygold}
Jede hinreichend große gerade ganze Zahl lässt sich als Summe zweier Werte
der Form $f(m)$ und $f(n)$ mit prim schreiben, für ein gegebenes
primdarstellendes Polynom $f$.
\end{vermutung}

Konkrete Fragen umfassen:
\begin{itemize}[nosep]
  \item Ist jede ganze Zahl als $p + n^2$ mit $p$ prim darstellbar?
  \item Gibt es unendlich viele Primzahlen der Form $n^2 + 1$?
        (Landaus viertes Problem, 1912)
\end{itemize}

\begin{bemerkung}
Landaus viertes Problem (1912) fragt explizit nach unendlich vielen Primzahlen
der Form $n^2 + 1$. Es ist ein Spezialfall der Bunyakovsky-Vermutung und
vollständig offen.
\end{bemerkung}

% ============================================================
\section{Bekannte Ergebnisse in Richtung Bunyakovsky}
% ============================================================

\begin{satz}[Iwaniec 1978]
Das Polynom $f(n) = n^2 + 1$ nimmt unendlich viele Werte an, die entweder
prim oder Produkt von genau zwei Primzahlen (fast-prim, $P_2$) sind.
\end{satz}

\begin{bemerkung}
Iwaniecs Ergebnis verwendet das Sieb von Eratosthenes in verfeinerter Form
(Brun--Titchmarsh--Iwaniec-Sieb).
Es zeigt, dass $n^2 + 1$ nicht stets zusammengesetzt sein kann, beweist aber
nicht die Unendlichkeit der Primwerte, weil das Sieb das ``Paritätsproblem''
nicht löst.
\end{bemerkung}

\begin{satz}[Linnik für lineare Polynome]
Die kleinste Primzahl in der Progression $a \pmod b$ ist $\ll b^L$
für eine absolute Konstante $L$. Aktuell bestes Ergebnis: $L = 5$
(Xylouris 2011).
\end{satz}

% ============================================================
\section{Übersicht bekannter Fälle und offener Probleme}
% ============================================================

\begin{center}
\begin{tabular}{lll}
\toprule
Polynom & Status & Referenz \\
\midrule
$an + b$, $\gcd(a,b)=1$ & Bewiesen (unendlich viele Primzahlen) & Dirichlet 1837 \\
$n^2 + 1$ & Offen (vermutet unendlich) & Bunyakovsky 1857 \\
$n^2 + n + 41$ & Offen (vermutet unendlich) & Bunyakovsky 1857 \\
$n^4 + 1$ & Offen & Bunyakovsky 1857 \\
$n^2 + 1$ ist $P_2$ & Bewiesen (unendlich viele Fastprimzahlen) & Iwaniec 1978 \\
\bottomrule
\end{tabular}
\end{center}

\noindent\textbf{Offene Fragen:}
\begin{enumerate}[label=(\arabic*)]
  \item Beweise die Bunyakovsky-Vermutung für ein einzelnes irreduzibles
        Polynom vom Grad $\geq 2$.
  \item Beweise Hypothese H für irgendein Paar von Polynomen.
  \item Löse Landaus viertes Problem: unendlich viele Primzahlen der Form $n^2+1$.
  \item Verbessere die Bateman--Horn-Konstante $C(f)$ durch unbedingte Methoden.
\end{enumerate}

% ============================================================
\section{Schluss}
% ============================================================

Die Bunyakovsky-Vermutung ist die natürliche Verallgemeinerung von Dirichlets Satz
auf Polynome vom Grad $\geq 2$.
Ihr einfachster Spezialfall --- dass $n^2 + 1$ unendlich oft prim ist ---
hat dem Beweis seit über 165 Jahren widerstanden.
Die Vermutung wird durch massive Rechenevidenz, die quantitative
Bateman--Horn-Vorhersage und Analogien mit bewiesenen Ergebnissen in
Funktionenkörpern gestützt.
Dennoch erscheint ein Beweis in weiter Ferne.

\begin{vermutung}[Wiederholt]
Jedes irreduzible $f \in \Z[x]$ mit positivem führenden Koeffizienten
und $\Delta(f) = 1$ stellt unendlich viele Primzahlen dar.
\end{vermutung}

% ============================================================
\begin{thebibliography}{99}

\bibitem{Bunyakovsky1857}
V.\ Bunyakovsky,
\textit{Nouveaux théorèmes relatifs à la distinction des nombres premiers},
Mem.\ Acad.\ Sci.\ St.\,Pétersbourg (6) \textbf{6} (1857), 305--329.

\bibitem{Dirichlet1837}
P.\ G.\ L.\ Dirichlet,
\textit{Beweis des Satzes, daß jede unbegrenzte arithmetische Progression
unendlich viele Primzahlen enthält},
Abh.\ Kön.\ Preuss.\ Akad.\ Wiss.\ (1837), 45--81.

\bibitem{SchinzelSierpinski1958}
A.\ Schinzel und W.\ Sierpiński,
\textit{Sur certaines hypothèses concernant les nombres premiers},
Acta Arith.\ \textbf{4} (1958), 185--208.

\bibitem{BatemanHorn1962}
P.\ T.\ Bateman und R.\ A.\ Horn,
\textit{A heuristic asymptotic formula concerning the distribution of prime numbers},
Math.\ Comp.\ \textbf{16} (1962), 363--367.

\bibitem{Iwaniec1978}
H.\ Iwaniec,
\textit{Almost-primes represented by quadratic polynomials},
Invent.\ Math.\ \textbf{47} (1978), 171--188.

\bibitem{Xylouris2011}
T.\ Xylouris,
\textit{On the least prime in an arithmetic progression},
Acta Arith.\ \textbf{150} (2011), 65--91.

\bibitem{CrandallPomerance2005}
R.\ Crandall und C.\ Pomerance,
\textit{Prime Numbers: A Computational Perspective},
2.\ Aufl., Springer, 2005.

\bibitem{Ono2004}
K.\ Ono,
\textit{The Web of Modularity},
CBMS \textbf{102}, AMS, 2004.

\bibitem{Granville1995}
A.\ Granville,
\textit{Harald Cramér and the distribution of prime numbers},
Scand.\ Actuar.\ J.\ \textbf{1995}, 12--28.

\end{thebibliography}

\end{document}
